% ======================================================================
%  Mecánica Clásica usando IA — Sesión 13, Tema 3
%  Paréntesis de Poisson
%  Codificación: UTF-8.   Compilar con pdflatex (dos veces, por el índice).
%  Duración estimada: 90 min de exposición + actividades IA fuera de clase
%  Referencia: M. Cosenza, Mecánica Clásica (2024), Secs. 6.4 y 6.5
%              (ecs. 6.75-6.103 y 6.173-6.177).
%  Sesión previa: transformaciones canónicas y funciones generadoras.
%  Sesión siguiente: Problemas de dinámica hamiltoniana.
% ======================================================================
\documentclass[spanish]{beamer}
\usepackage[utf8]{inputenc}
\usepackage[spanish]{babel}      % define \sen y usa coma decimal
\usepackage{amsmath,amssymb,mathtools}
\usepackage{graphicx}
\usepackage{textpos}

\usetheme{Madrid}
\usecolortheme{beaver}
\setbeamertemplate{navigation symbols}{}

% ---------------------------------------------------------------------
%  Logo (cambie sólo esta ruta; si el archivo no existe, se omite)
% ---------------------------------------------------------------------
\newcommand{\logopath}{/Users/luisnunez/Dropbox/MisDocumentos/UIS/UISImagenInstitucional/UISLOGO.png}
\newcommand{\logoUIS}{\IfFileExists{\logopath}{\includegraphics[height=0.4in]{\logopath}}{}}
\addtobeamertemplate{frametitle}{}{%
  \begin{textblock*}{100mm}(.85\textwidth,-1cm)\logoUIS\end{textblock*}}

% ---------------------------------------------------------------------
%  Notación unificada del curso
%  L caligráfica: lagrangiano;  H caligráfica: hamiltoniano;
%  L (recta): momento angular;  sen: babel-spanish
% ---------------------------------------------------------------------
\providecommand{\sen}{\operatorname{sen}}
\newcommand{\Lag}{\mathcal{L}}
\newcommand{\Ham}{\mathcal{H}}
\newcommand{\pd}[2]{\frac{\partial #1}{\partial #2}}
\newcommand{\pb}[2]{\left\{#1,#2\right\}}   % Cosenza escribe [f,g]
\newcommand{\J}{\mathsf{J}}                  % matriz simpléctica

% ---------------------------------------------------------------------
\title[Paréntesis de Poisson]{\textbf{Paréntesis de Poisson}}
\subtitle{La estructura algebraica de la dinámica hamiltoniana}
\author[L.A. Núñez]{\textbf{Luis A. Núñez}}
\institute[UIS]{\textit{Escuela de Física, Facultad de Ciencias,} \\
  \textit{Universidad Industrial de Santander, Santander, Colombia} \\
  \logoUIS}
\date{\today}

\begin{document}

\begin{frame}
  \titlepage
\end{frame}

\begin{frame}{Agenda}
  \tableofcontents
\end{frame}

% ======================================================================
\section{Objetivos y convenciones}
% ======================================================================

\begin{frame}{Objetivos de la sesión}
\small
Al finalizar, el estudiante podrá:
\begin{enumerate}
  \item<1-> \textbf{Derivar} la evolución temporal de cualquier función del espacio
        de fase como $df/dt = \pb{f}{\Ham} + \partial f/\partial t$.
  \item<2-> \textbf{Usar} las propiedades algebraicas del paréntesis (bilinealidad,
        antisimetría, Leibniz, Jacobi) para calcular sin derivar término a término.
  \item<3-> \textbf{Identificar} constantes de movimiento y \textbf{obtener} otras con
        el teorema de Poisson.
  \item<4-> \textbf{Interpretar} una función $G$ como generadora de una transformación
        canónica infinitesimal y \textbf{relacionar} simetrías con conservación.
  \item<5-> \textbf{Determinar} si una transformación es canónica mediante los
        paréntesis fundamentales.
\end{enumerate}
\vspace{0.2cm}
\onslide<6->{\textbf{Prerrequisitos:} ecuaciones de Hamilton, coordenadas cíclicas,
  teorema de Noether, transformaciones canónicas y funciones generadoras
  (Cosenza, Cap.~6).}
\end{frame}

\begin{frame}{Notación y convenciones del curso}
\begin{itemize}
  \item Sistema con $s$ grados de libertad; espacio de fase $(q_i,p_i)$,
        $i=1,\ldots,s$; hamiltoniano $\Ham(q_i,p_i,t)$.
  \item Vector de estado $\mathbf z = (q_1,\ldots,q_s,p_1,\ldots,p_s)^{\mathsf T}$ y
        matriz simpléctica
  \[
    \J = \begin{pmatrix} 0 & \mathsf I_s \\ -\mathsf I_s & 0 \end{pmatrix},
    \qquad \J^{\mathsf T} = -\J = \J^{-1}.
  \]
  \item Paréntesis de Poisson: $\pb{f}{g}$. Cosenza escribe $[f,g]$, con el
        \textbf{mismo signo}. Reservamos $[\hat A,\hat B]$ para el conmutador cuántico.
  \item $\Lag$, $\Ham$ caligráficas; $L$ (recta) es momento angular
        (Cosenza usa $l$). $\Ham'$ es el nuevo hamiltoniano tras una transformación
        canónica.
\end{itemize}
\end{frame}

\begin{frame}{Política de uso de IA en esta sesión}
\begin{itemize}
  \item<1-> \textbf{IA prohibida:} las derivaciones en clase y los chequeos
        conceptuales. Calcular paréntesis a mano es la destreza que se entrena.
  \item<2-> \textbf{IA con análisis crítico requerido:} las Actividades 1--4.
        La IA produce una respuesta; usted la audita con contraejemplos, casos
        límite y verificación simbólica.
  \item<3-> \textbf{IA asistida:} los problemas \emph{Para la discusión}. Puede delegar
        álgebra verificable y depuración de código; conserva el planteamiento y la
        interpretación.
  \item<4-> En todos los casos se entrega la \textbf{bitácora 4D}: prompts, respuesta,
        qué se aceptó/rechazó/corrigió, cómo se verificó y razonamiento final.
\end{itemize}
\end{frame}

% ======================================================================
\section{Definición}
% ======================================================================

\begin{frame}{Motivación: ¿se conserva $f$ sin resolver el movimiento?}
Sea $f(q_i,p_i,t)$ una función en el espacio de fase. A lo largo de una trayectoria,
\[
  \frac{df}{dt} = \sum_{i=1}^{s}\left( \pd{f}{q_i}\dot q_i + \pd{f}{p_i}\dot p_i \right)
                + \pd{f}{t} .
\]
\pause
Las ecuaciones de Hamilton $\dot q_i = \partial\Ham/\partial p_i$,
$\dot p_i = -\partial\Ham/\partial q_i$ eliminan las velocidades:
\[
  \frac{df}{dt} = \sum_{i=1}^{s}\left( \pd{f}{q_i}\pd{\Ham}{p_i}
                - \pd{f}{p_i}\pd{\Ham}{q_i} \right) + \pd{f}{t} .
\]
\pause
\begin{block}{Idea}
La variación de $f$ queda escrita sólo con derivadas de $f$ y de $\Ham$, evaluadas
en el punto del espacio de fase. No hace falta conocer $q_i(t)$, $p_i(t)$.
\end{block}
\end{frame}

\begin{frame}{Definición y criterio de conservación}
\begin{block}{Paréntesis de Poisson (Cosenza, ec.~6.80)}
Para $f(q_i,p_i,t)$ y $g(q_i,p_i,t)$:
\[
  \pb{f}{g} \equiv \sum_{i=1}^{s}\left( \pd{f}{q_i}\pd{g}{p_i}
              - \pd{f}{p_i}\pd{g}{q_i} \right)
            = (\nabla_{\mathbf z} f)^{\mathsf T}\,\J\,\nabla_{\mathbf z} g .
\]
\end{block}
\begin{itemize}
  \item<2-> Con $g=\Ham$:
  \(
    \dfrac{df}{dt} = \pb{f}{\Ham} + \pd{f}{t}.
  \)
  \item<3-> $f$ es constante de movimiento $\iff \pb{f}{\Ham} + \partial f/\partial t = 0$.
  \item<4-> Si además $\partial f/\partial t = 0$: \ $f$ se conserva $\iff \pb{f}{\Ham}=0$.
  \item<5-> Caso $f=\Ham$: $d\Ham/dt = \partial\Ham/\partial t$, pues
        $\pb{\Ham}{\Ham}=0$.
  \item<6-> \alert{Cuidado:} una $f$ con dependencia explícita de $t$ también puede
        conservarse (ver Para la discusión, Prob.~6.17).
\end{itemize}
\end{frame}

% ======================================================================
\section{Propiedades algebraicas}
% ======================================================================

\begin{frame}{Propiedades algebraicas}
Para funciones $f,g,h$ en el espacio de fase y constantes $a,b$:
\begin{enumerate}
  \item<1-> \textbf{Antisimetría:} $\pb{f}{g} = -\pb{g}{f}$; \ en particular $\pb{f}{f}=0$.
  \item<2-> \textbf{Bilinealidad:} $\pb{af_1+bf_2}{g} = a\pb{f_1}{g} + b\pb{f_2}{g}$
        (y lo mismo en el segundo argumento). Si $c$ es constante, $\pb{f}{c}=0$.
  \item<3-> \textbf{Regla de Leibniz:} $\pb{f_1f_2}{g} = f_1\pb{f_2}{g} + f_2\pb{f_1}{g}$.
  \item<4-> \textbf{Identidad de Jacobi:}
        $\pb{f}{\pb{g}{h}} + \pb{g}{\pb{h}{f}} + \pb{h}{\pb{f}{g}} = 0$.
  \item<5-> \textbf{Derivada temporal:}
        $\partial_t\pb{f}{g} = \pb{\partial_t f}{g} + \pb{f}{\partial_t g}$.
\end{enumerate}
\onslide<6->{
\begin{block}{Estructura}
1, 2 y 4 hacen de las funciones en el espacio de fase un \textbf{álgebra de Lie};
con 3, un \textbf{álgebra de Poisson}: $\pb{\cdot}{g}$ actúa como una derivada.
\end{block}}
\end{frame}

\begin{frame}{El paréntesis no es asociativo}
\begin{itemize}
  \item<1-> La regla de Leibniz \emph{no} expresa asociatividad: dice que
        $\pb{\cdot}{g}$ es una derivación. La asociatividad sería
        $\pb{f}{\pb{g}{h}} = \pb{\pb{f}{g}}{h}$, y es falsa en general.
  \item<2-> Contraejemplo con un grado de libertad:
  \[
    \pb{q}{\pb{q}{p^2}} = \pb{q}{2p} = 2,
    \qquad
    \pb{\pb{q}{q}}{p^2} = \pb{0}{p^2} = 0 .
  \]
  \item<3-> La identidad de Jacobi es lo que sustituye a la asociatividad. Su
        consecuencia más útil es el teorema de Poisson.
\end{itemize}
\end{frame}

\begin{frame}{Paréntesis con las coordenadas: paréntesis fundamentales}
\begin{itemize}
  \item<1-> Como $q_i$ y $p_i$ son variables independientes,
        $\partial q_i/\partial q_k = \delta_{ik}$ y $\partial q_i/\partial p_k = 0$:
  \[
    \pb{q_i}{f} = \sum_{k=1}^{s}\delta_{ik}\pd{f}{p_k} = \pd{f}{p_i},
    \qquad
    \pb{p_i}{f} = -\sum_{k=1}^{s}\delta_{ik}\pd{f}{q_k} = -\pd{f}{q_i}.
  \]
  \item<2-> Con $f = q_j$ o $f = p_j$ (\textbf{paréntesis fundamentales}):
  \[
    \pb{q_i}{q_j} = 0, \qquad \pb{p_i}{p_j} = 0, \qquad \pb{q_i}{p_j} = \delta_{ij} .
  \]
  \item<3-> Ecuaciones de Hamilton en forma simétrica:
  \[
    \dot q_i = \pb{q_i}{\Ham}, \qquad \dot p_i = \pb{p_i}{\Ham}
    \qquad\Longleftrightarrow\qquad
    \dot{\mathbf z} = \pb{\mathbf z}{\Ham} = \J\,\nabla_{\mathbf z}\Ham .
  \]
\end{itemize}
\end{frame}

% ======================================================================
\section{Momento lineal y angular}
% ======================================================================

\begin{frame}{Ejemplo 1: momento lineal y traslaciones}
Con $r = (x^2+y^2+z^2)^{1/2}$ y $\pb{r}{\mathbf p} \equiv \sum_j \pb{r}{p_j}\,\hat{\mathbf e}_j$:
\begin{itemize}
  \item<2-> Por la regla anterior, $\pb{f}{p_x} = -\pb{p_x}{f} = \partial f/\partial x$:
  \[
    \pb{r}{p_x} = \pd{r}{x} = \frac{x}{r},\qquad
    \pb{r}{p_y} = \frac{y}{r},\qquad
    \pb{r}{p_z} = \frac{z}{r},
  \]
  \[
    \Rightarrow\quad \pb{r}{\mathbf p} = \frac{\mathbf r}{r} = \hat{\mathbf r} .
  \]
  \item<3-> En general $\pb{f}{\mathbf p} = \nabla f$: el paréntesis con $\mathbf p$
        mide cómo cambia $f$ bajo una \textbf{traslación}.
  \item<4-> Por eso, si $\Ham$ es invariante bajo traslaciones,
        $\pb{\mathbf p}{\Ham} = -\nabla\Ham = 0$ y el momento total se conserva.
\end{itemize}
\end{frame}

\begin{frame}{Ejemplo 2: el álgebra del momento angular}
Con $L_x = yp_z - zp_y$, $L_y = zp_x - xp_z$, $L_z = xp_y - yp_x$:
\begin{itemize}
  \item<2-> Usando $\pb{p_i}{f} = -\partial f/\partial q_i$:
  \(
    \pb{p_y}{L_x} = -\partial L_x/\partial y = -p_z,
  \)
  \(
    \pb{p_x}{L_x} = 0,
  \)
  \(
    \pb{p_z}{L_y} = -p_x .
  \)
  \item<3-> Usando bilinealidad, Leibniz y los fundamentales (sólo sobreviven los
        términos con $z$ y $p_z$):
  \[
    \pb{L_x}{L_y} = \pb{yp_z - zp_y}{zp_x - xp_z}
                  = y\,p_x\pb{p_z}{z} + x\,p_y\pb{z}{p_z},
  \]
  \[
    \pb{L_x}{L_y} = -yp_x + xp_y = L_z .
  \]
  \item<4-> En forma compacta, con $\epsilon_{ijk}$ el símbolo de Levi-Civita:
  \[
    \pb{p_i}{L_j} = \epsilon_{ijk}\,p_k, \qquad
    \pb{L_i}{L_j} = \epsilon_{ijk}\,L_k, \qquad
    \pb{L^2}{L_i} = 0 .
  \]
\end{itemize}
\end{frame}

\begin{frame}{Consecuencias físicas del álgebra de $\mathbf L$}
\begin{itemize}
  \item<1-> Si $L_x$ y $L_y$ se conservan, entonces $L_z = \pb{L_x}{L_y}$ también
        (teorema de Poisson, próxima sección).
  \item<2-> Como $\pb{L_x}{L_y}\neq 0$, dos componentes de $\mathbf L$ \textbf{no}
        pueden ser simultáneamente momentos canónicos de un mismo conjunto
        $(Q_i,P_i)$.
  \item<3-> Fuerza central: $\Ham$, $L^2$ y $L_z$ son tres constantes con paréntesis
        mutuos nulos: el problema de $s=3$ es integrable (próxima sección).
  \item<4-> Componentes de un vector: $\pb{p_i}{L_j} = \epsilon_{ijk}p_k$ dice que
        $\mathbf p$ rota como un vector. Escalares como $p^2$, $r^2$ o
        $\mathbf r\cdot\mathbf p$ cumplen $\pb{f}{L_j} = 0$.
\end{itemize}
\end{frame}

\begin{frame}{Correspondencia con la mecánica cuántica}
\begin{itemize}
  \item<1-> En mecánica cuántica los observables son operadores y su conmutador es
        $[\hat A,\hat B] = \hat A\hat B - \hat B\hat A$.
  \item<2-> Regla de correspondencia de Dirac:
  \[
    \pb{f}{g} \;\longrightarrow\; \frac{1}{i\hbar}\,[\hat f,\hat g] .
  \]
  \item<3-> Así, $\pb{q_i}{p_j} = \delta_{ij}$ pasa a
        $[\hat q_i,\hat p_j] = i\hbar\,\delta_{ij}$, y
        $\pb{L_i}{L_j} = \epsilon_{ijk}L_k$ pasa a
        $[\hat L_i,\hat L_j] = i\hbar\,\epsilon_{ijk}\hat L_k$
        ($\hbar$: constante de Planck reducida).
  \item<4-> \alert{Límite:} la correspondencia no puede ser exacta para todas las
        funciones (orden de factores: ¿$qp$ es $\hat q\hat p$, $\hat p\hat q$ o su
        promedio?). Se preserva la estructura de las relaciones fundamentales,
        no cada fórmula clásica.
\end{itemize}
\end{frame}

% ======================================================================
\section{Teorema de Poisson e integrabilidad}
% ======================================================================

\begin{frame}{Teorema de Poisson}
\begin{block}{Enunciado (Cosenza, ecs.~6.101--6.102)}
Si $f$ y $g$ son constantes de movimiento, entonces $\pb{f}{g}$ también lo es.
\end{block}
\begin{itemize}
  \item<2-> Por la identidad de Jacobi,
        $\pb{\pb{f}{g}}{\Ham} = \pb{\pb{f}{\Ham}}{g} + \pb{f}{\pb{g}{\Ham}}$.
  \item<3-> Añadiendo $\partial_t\pb{f}{g}$ (propiedad 5):
  \[
    \frac{d}{dt}\pb{f}{g} = \pb{\frac{df}{dt}}{g} + \pb{f}{\frac{dg}{dt}} = 0 .
  \]
  \item<4-> \alert{Cuidado:} la nueva constante puede ser trivial (un número) o una
        función de las ya conocidas. Ejemplo: $\pb{q}{p}=1$.
        El teorema no garantiza una integral \emph{independiente}.
\end{itemize}
\end{frame}

\begin{frame}{Integrabilidad de Liouville}
\begin{block}{Condición (Cosenza, ec.~6.103)}
Un sistema de $s$ grados de libertad es integrable si existen $s$ constantes de
movimiento independientes $I_k(q,p)$ en \textbf{involución}:
$\pb{I_k}{I_j} = 0$ para todo $k,j$.
\end{block}
\begin{itemize}
  \item<2-> En sistemas conservativos, $\Ham$ es una de ellas.
  \item<3-> Fuerza central ($s=3$): $\{\Ham, L^2, L_z\}$ están en involución.
        Las tres componentes $L_x, L_y, L_z$ también se conservan, pero no están en
        involución entre sí.
  \item<4-> En cambio, el péndulo con soporte sobre la parábola ($s=2$) sólo
        tiene $\Ham$. Sus secciones de Poincaré muestran islas y caos
        (video de la sesión de problemas).
  \item<5-> Conexión: en involución significa ``observables que conmutan''. Las $I_k$
        serán los nuevos momentos en Hamilton--Jacobi y acción--ángulo.
\end{itemize}
\end{frame}

% ======================================================================
\section{Generadores de transformaciones}
% ======================================================================

\begin{frame}{Transformaciones canónicas infinitesimales}
\begin{itemize}
  \item<1-> Tome la función generadora cercana a la identidad, con $\varepsilon\ll 1$:
  \[
    F_2(q,P) = \sum_i q_iP_i + \varepsilon\,G(q,P)
    \;\Rightarrow\;
  \]
  \[
    p_i = P_i + \varepsilon\pd{G}{q_i},\qquad Q_i = q_i + \varepsilon\pd{G}{P_i}.
  \]
  \item<2-> A primer orden en $\varepsilon$ ($G(q,P)\approx G(q,p)$):
  \[
    \delta q_i = \varepsilon\,\pb{q_i}{G}, \qquad
    \delta p_i = \varepsilon\,\pb{p_i}{G}, \qquad
    \delta f = \varepsilon\,\pb{f}{G} .
  \]
  \item<3-> $G$ es el \textbf{generador} de la transformación. Ejemplos:
  \begin{itemize}
    \item $G=p_x$: \ $\delta x = \varepsilon$ (traslación).
    \item $G=L_z$: \ $\delta x = -\varepsilon y$, $\delta y = \varepsilon x$
          (rotación de ángulo $\varepsilon$ alrededor de $z$).
    \item $G=\Ham$ con $\varepsilon = dt$: \ $\delta\mathbf z = \dot{\mathbf z}\,dt$
          (la evolución temporal es una transformación canónica).
  \end{itemize}
\end{itemize}
\end{frame}

\begin{frame}{Simetrías y conservación: Noether en lenguaje hamiltoniano}
\begin{itemize}
  \item<1-> Bajo la transformación generada por $G$ (con $\partial G/\partial t = 0$):
  \[
    \delta\Ham = \varepsilon\,\pb{\Ham}{G} = -\varepsilon\,\pb{G}{\Ham}
               = -\varepsilon\,\frac{dG}{dt} .
  \]
  \item<2-> Por lo tanto
  \[
    \boxed{\ \Ham \text{ invariante bajo la transformación de } G
    \iff dG/dt = 0\ }
  \]
  \item<3-> Coordenada cíclica: si $q_k$ no aparece en $\Ham$, la traslación
        $\delta q_k=\varepsilon$ (generada por $p_k$) deja $\Ham$ invariante y
        $p_k$ se conserva.
  \item<4-> Una misma función tiene dos papeles: \textbf{cantidad observable}
        y \textbf{generador de una simetría}. Es la versión hamiltoniana del
        teorema de Noether visto en la formulación lagrangiana.
\end{itemize}
\end{frame}

% ======================================================================
\section{Paréntesis de Poisson y transformaciones canónicas}
% ======================================================================

\begin{frame}{Criterio de canonicidad e invariancia}
\begin{block}{Teorema (Cosenza, ecs.~6.174--6.177)}
$(q,p)\to(Q,P)$ es canónica si y sólo si, calculados en las variables $(q,p)$,
\[
  \pb{Q_i}{Q_j} = 0, \qquad \pb{P_i}{P_j} = 0, \qquad \pb{Q_i}{P_j} = \delta_{ij} .
\]
Además, $\pb{f}{g}_{q,p} = \pb{f}{g}_{Q,P}$ para cualquier par de funciones.
\end{block}
\begin{itemize}
  \item<2-> Forma matricial: con $\mathsf M = \partial(Q,P)/\partial(q,p)$, la condición
        es $\mathsf M\,\J\,\mathsf M^{\mathsf T} = \J$, lo que implica
        $\det\mathsf M = 1$: se conserva el volumen (Liouville).
  \item<3-> Número de condiciones independientes: $s(2s-1)$.
        Para $s=1$, sólo $\pb{Q}{P}=1$; para $s=2$, \alert{seis}.
  \item<4-> La invariancia permite calcular paréntesis en las variables más cómodas:
        por eso son ``la'' estructura de la mecánica, no un artificio de coordenadas.
\end{itemize}
\end{frame}

\begin{frame}{Ejemplo 3: un grado de libertad}
Sea $Q = q + e^{p}$, $P = p$.
\begin{itemize}
  \item<2-> $\pb{Q}{Q} = \pb{P}{P} = 0$ se cumplen siempre por antisimetría:
        no hace falta calcularlos.
  \item<3-> La única condición es
  \[
    \pb{Q}{P} = \pd{Q}{q}\pd{P}{p} - \pd{Q}{p}\pd{P}{q} = (1)(1) - (e^{p})(0) = 1 .
  \]
  \item<4-> Más rápido, con bilinealidad: $\pb{q + e^{p}}{p} = \pb{q}{p} + \pb{e^{p}}{p}
        = 1 + 0$. \quad \textbf{Es canónica.}
  \item<5-> En la sesión de problemas se usa esta transformación con
        $\Ham = q + t\,e^{p}$, y se construye su función generadora $F_2 = qP + e^{P}$.
\end{itemize}
\end{frame}

\begin{frame}{Ejemplo 4: dos grados de libertad}
Sea $Q_1 = q_1$, \ $Q_2 = p_2$, \ $P_1 = p_1 - 2p_2$, \ $P_2 = -2q_1 - q_2$.
Con bilinealidad y los paréntesis fundamentales:
\begin{align*}
  \pb{Q_1}{P_1} &= \pb{q_1}{p_1} = 1, &
  \pb{Q_2}{P_2} &= -\pb{p_2}{q_2} = 1, \\
  \pb{Q_1}{P_2} &= -\pb{q_1}{q_2} = 0, &
  \pb{Q_2}{P_1} &= -2\pb{p_2}{p_2} = 0, \\
  \pb{Q_1}{Q_2} &= \pb{q_1}{p_2} = 0 .
\end{align*}
\onslide<2->{El paréntesis entre momentos (sólo sobreviven los términos $p_iq_i$):
\[
  \pb{P_1}{P_2} = \pb{p_1}{-2q_1} + \pb{-2p_2}{-q_2} = 2 - 2 = 0 .
\]}
\onslide<3->{Las \textbf{seis} condiciones se cumplen: \textbf{es canónica}.
\alert{Las dos últimas también deben verificarse}; omitirlas es un error frecuente
(Actividad 3).}
\end{frame}

% ======================================================================
\section{Recapitulando}
% ======================================================================

\begin{frame}{Chequeos conceptuales (sin IA, 10 minutos)}
\begin{enumerate}
  \item<1-> Si $\pb{f}{\Ham}=0$, ¿se conserva $f$? ¿Qué condición falta?
  \item<2-> Sin calcular: ¿cuánto vale $\pb{f(\mathbf r)}{L_z}$ si $f$ sólo depende
        de $r=|\mathbf r|$? ¿Por qué?
  \item<3-> ¿Puede $L_x$ ser una coordenada $Q_1$ y $L_y$ su momento conjugado $P_1$?
  \item<4-> Con $s=2$, alguien verifica $\pb{Q_i}{P_j}=\delta_{ij}$ y concluye que la
        transformación es canónica. ¿Qué le falta?
  \item<5-> ¿Qué transformación genera $G = \mathbf r\cdot\mathbf p$? ¿Cuándo es una
        simetría?
\end{enumerate}
\end{frame}
% NOTA DOCENTE (respuestas breves):
% 1. Hace falta df/dt = 0, i.e., que f no dependa explícitamente de t
%    (o que {f,H} + df/dt parcial = 0).
% 2. Cero: f(r) es invariante bajo rotaciones alrededor de z.
% 3. No: {L_x, L_y} = L_z != 1.
% 4. {Q_1,Q_2} = 0 y {P_1,P_2} = 0.
% 5. Dilataciones: delta r = eps r, delta p = -eps p. Es simetría si H es
%    invariante de escala (no lo es en general); ver Prob. 6.17.

\begin{frame}{Recapitulando}
\begin{itemize}
  \item<1-> $\dfrac{df}{dt} = \pb{f}{\Ham} + \pd{f}{t}$: la dinámica entera cabe en
        $\dot{\mathbf z} = \J\nabla\Ham$.
  \item<2-> Bilinealidad, antisimetría, Leibniz y Jacobi permiten calcular paréntesis
        sin derivar término a término.
  \item<3-> Teorema de Poisson: de dos constantes se obtiene una tercera (que puede
        ser trivial). Integrabilidad: $s$ constantes en involución.
  \item<4-> Cada función $G$ genera una transformación canónica infinitesimal
        $\delta f = \varepsilon\pb{f}{G}$. Simetría de $\Ham$ $\iff$ $G$ conservada.
  \item<5-> Canónica $\iff$ paréntesis fundamentales preservados
        ($s(2s-1)$ condiciones) $\iff \mathsf M\J\mathsf M^{\mathsf T} = \J$.
\end{itemize}
\vspace{0.2cm}
\onslide<6->{\textbf{Próximas sesiones:} problemas de dinámica hamiltoniana; ecuación de
  Hamilton--Jacobi y variables de acción--ángulo.}
\end{frame}

% ======================================================================
\section{Actividades con IA (modelo 4D)}
% ======================================================================

\begin{frame}{Actividad 1 (Discernimiento): auditar un resumen}
Un asistente de IA produce este resumen de propiedades:
\begin{enumerate}
  \item[(a)] $\pb{f_1f_2}{g} = f_1\pb{f_2}{g} + f_2\pb{f_1}{g}$ \ (propiedad asociativa).
  \item[(b)] $\pb{f}{\pb{g}{\Ham}} + \pb{g}{\pb{h}{f}} + \pb{h}{\pb{f}{g}} = 0$ \ (Jacobi).
  \item[(c)] En mecánica cuántica, $\pb{q_i}{p_j} = i\hbar\,\delta_{ij}$.
  \item[(d)] Si $f$ no depende de $t$ y se conserva, entonces $\pb{f}{\Ham} = 0$.
\end{enumerate}
\begin{itemize}
  \item<2-> \textbf{Discernimiento:} clasifique cada afirmación (correcta, incorrecta,
        mal nombrada) y justifique con un contraejemplo o una corrección.
  \item<3-> \textbf{Descripción:} escriba un prompt que obligue al asistente a
        distinguir paréntesis clásico y conmutador cuántico.
  \item<4-> \textbf{Diligencia:} registre si el asistente detecta por sí mismo los
        errores cuando le pide revisar su propio resumen.
\end{itemize}
\end{frame}
% NOTA DOCENTE: (a) correcta pero mal nombrada (Leibniz, no asociatividad);
% (b) error tipográfico: debe ser {g,h}, no {g,H}; (c) incorrecta: es el
% conmutador [q,p] = i hbar delta; (d) correcta. Estos cuatro errores estaban
% en la versión anterior de estas diapositivas.

\begin{frame}[fragile]{Actividad 2 (Delegación): paréntesis en SymPy}
Delegue el álgebra y conserve la interpretación:
\footnotesize
\begin{verbatim}
import sympy as sp
x, y, z, px, py, pz = sp.symbols('x y z p_x p_y p_z', real=True)
q, p = [x, y, z], [px, py, pz]

def pb(f, g):
    return sp.simplify(sum(sp.diff(f, q[i])*sp.diff(g, p[i])
                           - sp.diff(f, p[i])*sp.diff(g, q[i])
                           for i in range(3)))

Lx, Ly, Lz = y*pz - z*py, z*px - x*pz, x*py - y*px
print(pb(Lx, Ly) - Lz)        # debe dar 0
\end{verbatim}
\normalsize
\begin{itemize}
  \item Extienda: verifique $\pb{L^2}{L_i}=0$ y la identidad de Jacobi para tres
        polinomios de su elección.
  \item Discernimiento: ¿por qué verificar Jacobi con \emph{tres} ejemplos no la
        demuestra? ¿Qué sí sería una demostración?
\end{itemize}
\end{frame}

\begin{frame}{Actividad 3 (Discernimiento): ¿es canónica?}
Pida a un asistente que decida si la siguiente transformación ($s=2$) es canónica:
\[
  Q_1 = q_1, \qquad Q_2 = q_2, \qquad P_1 = p_1 + q_2, \qquad P_2 = p_2 .
\]
\begin{itemize}
  \item<2-> Calcule usted los seis paréntesis. ¿Cuáles se cumplen y cuál falla?
  \item<3-> Si el asistente sólo verifica $\pb{Q_i}{P_j} = \delta_{ij}$, documente el
        error y explique por qué no basta.
  \item<4-> \textbf{Extensión:} modifique $P_2$ para que la transformación sea canónica
        y encuentre una función $F_2(q,P)$ que la genere.
        Interprete el resultado como un ``potencial vector'' en el espacio de
        configuración.
\end{itemize}
\end{frame}
% NOTA DOCENTE: {P_1,P_2} = {q_2,p_2} = 1 != 0: no canónica, aunque las
% cuatro {Q_i,P_j} y {Q_1,Q_2} se cumplen. Corrección: P_2 = p_2 + q_1,
% con F_2 = q_1 P_1 + q_2 P_2 - q_1 q_2 (p_i = P_i - dLambda/dq_i,
% Lambda = q_1 q_2): el corrimiento debe ser un gradiente.

\begin{frame}{Actividad 4 (Descripción): del paréntesis al conmutador}
\begin{itemize}
  \item<1-> Pida a un asistente que explique la regla
        $\pb{f}{g}\to [\hat f,\hat g]/(i\hbar)$ y que dé dos ejemplos.
  \item<2-> \textbf{Discernimiento:} compare $\pb{x}{L_z}$ con
        $[\hat x,\hat L_z]/(i\hbar)$ calculado con $[\hat x,\hat p_x] = i\hbar$.
        ¿Coinciden?
  \item<3-> Pregunte qué ocurre con el orden de factores en $qp$. Exija que el
        asistente distinga la estructura algebraica que se preserva de las
        fórmulas que no se trasladan término a término. Verifique la afirmación en una
        referencia de mecánica cuántica antes de aceptarla.
  \item<4-> \textbf{Diligencia:} anote en la bitácora qué afirmaciones pudo verificar y
        cuáles quedaron sin respaldo.
\end{itemize}
\end{frame}
% NOTA DOCENTE: {x,L_z} = dL_z/dp_x = -y; [x, L_z] = [x, x p_y - y p_x]
% = -y [x,p_x] = -i hbar y, de modo que [x,L_z]/(i hbar) = -y: coinciden.

\begin{frame}{Diligencia: bitácora 4D y declaración de uso de IA}
Cada actividad con IA se entrega con una bitácora que contiene:
\begin{enumerate}
  \item Los prompts usados (completos y en orden).
  \item Las partes relevantes de cada respuesta.
  \item Qué se aceptó, rechazó o corrigió, y por qué.
  \item Cómo se verificó: cálculo a mano de al menos un paréntesis, contraejemplos,
        verificación simbólica, consistencia con los paréntesis fundamentales.
  \item El razonamiento final, escrito por usted.
  \item Errores, omisiones o ``alucinaciones'' detectados.
\end{enumerate}
\vspace{0.2cm}
\textbf{Verificación de autoría:} cada estudiante calculará en el tablero (5 min) un
paréntesis de su bitácora con un cambio propuesto por el docente.
\end{frame}

% ======================================================================
\section{Para la discusión}
% ======================================================================

\begin{frame}{Para la discusión 1/3: fundamentos e integración}
\textbf{(Nivel 1)} Cosenza, Prob.~6.19: evalúe $\pb{\mathbf L}{\mathbf r\cdot\mathbf p}$,
$\pb{\mathbf p}{r^n}$ y $\pb{\mathbf p}{(\mathbf a\cdot\mathbf r)^n}$, con $\mathbf a$
constante.
\begin{itemize}
  \item IA: prohibida. Use bilinealidad y Leibniz antes que derivadas explícitas.
\end{itemize}
\vspace{0.2cm}
\textbf{(Nivel 2)} Cosenza, Prob.~6.14: péndulo esférico de longitud $a$ en $(\theta,\phi)$.
Escriba $\Ham$ y las ecuaciones de Hamilton y evalúe $\pb{L_x}{p_\phi}$.
\begin{itemize}
  \item Discernimiento: ¿qué cantidad física es $p_\phi$? ¿El resultado era previsible
        del álgebra de $\mathbf L$?
\end{itemize}
\vspace{0.2cm}
\textbf{(Nivel 2)} Cosenza, Prob.~6.17: $\Ham = |\mathbf p|^n - a|\mathbf r|^{-n}$.
¿Cuál de $f = \mathbf r\cdot\mathbf p/n - \Ham t$ y $g = |\mathbf p|$ es constante de
movimiento?
\end{frame}
% NOTA DOCENTE:
% 6.19: {L_i, r.p} = 0 (escalar); {p, r^n} = -n r^(n-2) r;
%       {p, (a.r)^n} = -n (a.r)^(n-1) a.
% 6.14: L_x = -sen(phi) p_theta - cot(theta) cos(phi) p_phi, p_phi = L_z;
%       {L_x, p_phi} = dL_x/dphi = -L_y.
% 6.17: {r.p, H} = n H => df/dt = H - H = 0: f se conserva (con t explícito).
%       {|p|, H} = -a n r^(-n-2)(r.p)/|p| != 0: g no se conserva.

\begin{frame}{Para la discusión 2/3: transferencia}
\textbf{(Nivel 3)} Cosenza, Prob.~6.18: vector de Laplace--Runge--Lenz para Kepler,
$\mathbf A = \mathbf p\times\mathbf L - mk\,\hat{\mathbf r}$, con
$\Ham = p^2/2m - k/r$.
\begin{enumerate}
  \item Calcule $\pb{A_i}{L_j}$ e interprete el resultado.
  \item Muestre que $\pb{A_i}{\Ham}=0$. ¿Qué simetría ``oculta'' revela?
  \item Calcule $\pb{A_i}{A_j}$. ¿Da el teorema de Poisson una constante nueva?
\end{enumerate}
\begin{itemize}
  \item IA asistida. \textbf{Delegación:} delegue el álgebra a SymPy
        (Actividad 2); conserve la interpretación. \textbf{Discernimiento:}
        cuente las constantes independientes. ¿Por qué las órbitas ligadas son
        cerradas?
\end{itemize}
\end{frame}
% NOTA DOCENTE: {A_i,L_j} = eps_ijk A_k (A es un vector); {A_i,H} = 0;
% {A_i,A_j} = -2 m H eps_ijk L_k: Poisson no da nada nuevo. Constantes:
% H, L (3), A (3) con A.L = 0 y A^2 = m^2k^2 + 2mH L^2 => 5 independientes
% (superintegrable): órbitas cerradas.

\begin{frame}{Para la discusión 3/3: investigación}
\textbf{(Nivel 4)} Integrable frente a caótico.
\begin{enumerate}
  \item Cosenza, Prob.~6.15 (red de Toda): verifique simbólicamente que
        $\pb{I}{\Ham}=0$ para la constante $I$ dada.
  \item Construya secciones de Poincaré para la red de Toda y compárelas con las del
        péndulo con soporte sobre la parábola (sesión de problemas). Describa las
        diferencias.
  \item Para el péndulo sobre la parábola, proponga dos candidatas a segunda integral y
        muestre que $\pb{f}{\Ham}\neq 0$. ¿Qué evidencia numérica sostiene que no
        existe ninguna?
\end{enumerate}
\textbf{IA:} asistida para código y álgebra; bitácora 4D obligatoria.
Entregables: verificación simbólica, código documentado, secciones de Poincaré
interpretadas y pruebas de convergencia.
\end{frame}
% NOTA DOCENTE: {I,H} = 0 se verifica con SymPy (expandir antes de simplificar).
% Las secciones de Toda muestran sólo curvas invariantes a cualquier energía;
% las del péndulo sobre la parábola (a=1, M=m=l=g=1) muestran islas a
% H - Hmin = 1,2 y mar caótico a 1,85.

\end{document}
